\chapter{The problem, and what it demands of the tooling}
\label{ch:problem}

\section{Physical target}

The driver is the characterisation of the effects and time-dependent dynamics
of the thermodiffusive instability in lean premixed hydrogen flames. These
flames spontaneously form cellular structures whose sizes and separations
change dramatically with fuel mixture and with the turbulence level of the
approach flow. The analysis lineage is that of Bremer
et~al.~\cite{bremer2010burning}; the present work is a rebuild of the
underlying geometric and topological machinery on modern AMReX.

Relevant configurations include a doubly periodic \emph{flamesheet} with inflow
and outflow in the third direction, and growing flame kernels, each in two and
three dimensions, and computed with different models for the transport physics
that drives the instability.

\section{What is not known}

The tooling exists because a prior question is open: \emph{what a cell is}.

There are chemistry arguments about when a local flame is self-sustaining, and
there is the distinct notion of \emph{back-supported} structure. Separately,
observation supplies a set of surrogates that have traditionally stood in for
the presence of a flame---fuel consumption rate, temperature above some
threshold, a flame intermediate, and others. Which surrogate to use, and at
what threshold, is not settled, and the answer may not be the same for all of
them.

A complication specific to these flames sharpens the difficulty. Because they
are far from unity Lewis number, the progress variables do not align: contours
of fuel, of temperature and of intermediates cross one another. Each may
therefore tell a different story in detail while supporting a similar global
picture---and whether that is so is itself one of the questions to be answered.

\section{Requirements that follow}

\begin{enumerate}[label=(R\arabic*),leftmargin=*]
\item\label{req:compare} \textbf{Compare, do not privilege.} The deliverable is
  not the topology of a chosen field but a framework in which candidate
  definitions are computed identically and compared directly. Adding a
  candidate surrogate should be a command-line argument, not a code change.
\item\label{req:robust} \textbf{Robustness to definition.} Any reported measure
  must come with a statement of how sensitive it is to the threshold and field
  that produced it. A number without that statement is not usable here.
\item\label{req:periodic} \textbf{Periodic correctness with arbitrary face
  placement.} Doubly and triply periodic domains are routine, and the location
  of the domain faces carries no physical meaning. No reported quantity may
  depend on where they happen to fall.
\item\label{req:amr} \textbf{Correctness across coarse--fine interfaces.} The
  data are adaptively refined; an analysis that is only approximately right at
  level transitions is not acceptable, because the structures of interest are
  exactly where the refinement is.
\item\label{req:parallel} \textbf{Parallel throughout.} Snapshots are large.
  Any step that gathers the whole problem to one rank is a design fault, not an
  optimisation opportunity.
\item\label{req:track} \textbf{Relate structures across snapshots.} The
  dynamics are the point, so features must be identifiable from one time to the
  next.
\item\label{req:sample} \textbf{Carry sampled data onto extracted geometry.}
  Surfaces must be able to carry an arbitrary vector of interpolated field
  values at their nodes; this is what makes conditional analysis possible.
\end{enumerate}

\Cref{req:robust} deserves emphasis because it determines the topological
target in \cref{ch:topology}: it is a request for \emph{persistence}, and
persistence is what a merge tree supplies.
